% 箭头自上而下，旁边标明交给下一步的模型、绘制数据或图像。
\begin{tikzpicture}[x=1mm,y=1mm,
  every node/.style={anchor=west,inner sep=0pt,outer sep=0pt,text=Ink,
    font=\sffamily\fontsize{10.5}{13}\selectfont},
  heading/.style={font=\sffamily\bfseries\fontsize{12}{15}\selectfont},
  handoff/.style={font=\sffamily\fontsize{10.3}{13}\selectfont,text=Muted},
  arrow/.style={draw=Ink,line width=.85pt,-{Stealth[length=2mm,width=1.75mm]}},
  stage/.style={rounded corners=1.5mm,line width=.65pt}]
  \definecolor{LoadBlue}{HTML}{305C80}
  \definecolor{DrawGreen}{HTML}{246659}
  \definecolor{PanelOchre}{HTML}{96502B}
  \draw[rounded corners=2mm,draw=Ink!72,fill=Paper,line width=.85pt]
    (0,0) rectangle (158,101);

  % 同一绿色的两步分别决定画什么、怎样交给后端。
  \foreach \bottom/\shade in {79/LoadBlue,54/DrawGreen,29/DrawGreen,4/PanelOchre}{
    \draw[stage,draw=\shade!80,fill=\shade!8] (4,\bottom) rectangle (154,\bottom+18);
    \draw[\shade!42,line width=.55pt] (36,\bottom+3.3)--(36,\bottom+14.7);
  }
  \node[heading,text=LoadBlue] at (8,88) {加载模型};
  \node at (40,91.5) {\texttt{AssetManager} 找到模型文件，加载器读取网格和材质。};
  \node at (40,84.5) {拟由 \texttt{RenderDevice} 创建顶点缓冲和砖墙贴图。};

  \node[heading,text=DrawGreen] at (8,63) {选择网格};
  \node at (40,66.5) {\texttt{SceneRenderer} 把模型、位置和编号交给 \texttt{Renderer3D}。};
  \node at (40,59.5) {\texttt{Renderer3D} 选出相机能看到的网格和对应的材质。};

  \node[heading,text=DrawGreen] at (8,38) {绘制场景};
  \node at (40,41.5) {拟由 \texttt{CommandList} 接收绘制命令，交给后端执行。};
  \node at (40,34.5) {后端将网格画进场景颜色图。现用 OpenGL，拟接入 Vulkan。};

  \node[heading,text=PanelOchre] at (8,13) {显示画面};
  \node at (40,16.5) {后处理调整曝光，写出面板输出图。};
  \node at (40,9.5) {编辑器用 ImGui 把这张图显示在场景面板里。};

  \foreach \high/\low/\what in {
    79/72/{加载好的模型，后续帧可以复用},
    54/47/{选定的网格、材质和物体位置},
    29/22/{画好的场景颜色图}}{
    \draw[arrow] (20,\high)--(20,\low);
    \node[handoff] at (24,{(\high+\low)/2}) {\what};
  }

  % 面板中的立方体是流程的输出示意。
  \begin{scope}[shift={(128,6.2)}]
    \filldraw[fill=Paper,draw=PanelOchre!80,line width=.5pt]
      (0,0) rectangle (22,13.6);
    \fill[PanelOchre!23] (.2,11.6) rectangle (21.8,13.3);
    \fill[PanelOchre!72] (5.4,8.5)--(10.6,6.5)--(10.6,1.2)--(5.4,3.2)--cycle;
    \fill[PanelOchre] (10.6,6.5)--(16.5,8.7)--(16.5,3.4)--(10.6,1.2)--cycle;
    \fill[PanelOchre!46] (5.4,8.5)--(11.3,10.7)--(16.5,8.7)--(10.6,6.5)--cycle;
    \draw[Paper!84!PanelOchre,line width=.35pt]
      (5.4,5.85)--(10.6,3.85)--(16.5,6.05)
      (8,7.5)--(8,4.85) (13.55,4.95)--(13.55,2.3);
  \end{scope}
\end{tikzpicture}
